\documentclass[11pt]{article}
\usepackage[margin=24mm]{geometry}
\usepackage{amsmath}
\usepackage{xcolor}
\usepackage{hyperref}
\hypersetup{hidelinks}
\setlength{\emergencystretch}{1em}
\newcommand{\response}[1]{\par\noindent\textcolor{blue!55!black}{\textbf{Response.} #1}\par}

\begin{document}

\begin{center}
{\Large Response to Technical Review dated 29 July 2026}\\[4pt]
{\bfseries GFNO-PINO: A Topology-Disjoint Benchmark Isolating Physics
Fine-Tuning, Topology Conditioning, and Spectral Filtering in AC-OPF Graph
Surrogates}
\end{center}

We thank the reviewer for the precise formulation, architecture, and
reproducibility audit. We corrected the manuscript and regenerated its tables,
figures, PDF, integrity manifest, and submission archive. We did not invent a
repository DOI or add unexecuted topology experiments. The resulting scope is
more conservative and matches the archived code and evidence.

\begin{enumerate}
\item \textbf{The N-1 connectivity equation omitted fixed transformers.}
\response{Equation~(3) now tests connectivity on the active full branch set
$\{a\in\mathcal B:s_a(z)=1\}$, with
$\mathcal B=\mathcal L\cup\mathcal X$. This matches the data generator, which
retains transformer and parallel-circuit edges while screening line outages.}

\item \textbf{The neural-operator/FNO terminology exceeded the evidence.}
\response{We now call the method a topology-conditioned, physics-informed
Chebyshev graph surrogate. We state that its localized polynomial layer is
structurally ChebNet, that each grid size has a separate checkpoint, and that
no mesh-refinement, discretization-transfer, or shared mixed-size model was
tested. ``GFNO-PINO'' is retained only as the archived experiment identifier.
The abstract, contributions, methodology, limitations, conclusion, model
docstrings, and README were revised consistently.}

\item \textbf{Branch angle limits were incorrectly described as predictor
inputs.}
\response{The formulation now separates predictor metadata
$\mathcal M_{\rm pred}$ from physics/evaluation metadata
$\mathcal M_{\rm phys}$. The former contains box limits, incidence, economics,
generator identity, and slack information. Branch angle limits and exact
$Y,Y_f,Y_t$ belong only to the latter. This matches the implemented
ten-channel edge tensor exactly.}

\item \textbf{Logistic heads cannot attain active endpoints.}
\response{The hard-layer section now states that finite logits map to the open
interval and therefore cannot reproduce an endpoint exactly. We evaluated all
20 GFNO-PINO checkpoints using a $10^{-4}$ p.u. near-bound definition. A new
table reports reference active fractions and active/inactive MAEs for
$V,P^g,Q^g$ in every case. The executable
\texttt{analyze\_active\_bounds.py} and its CSV are included. We discuss
projection/hard-sigmoid heads as a retraining-dependent alternative rather
than changing the archived experiment retrospectively.}

\item \textbf{Baseline training and checkpoint selection were underspecified.}
\response{A new protocol table gives fitting rows/topologies, validation
rows/topologies, training objective, checkpoint score, early-stopping status,
and exact parameter counts for every model and case. It discloses that the
fixed PINN uses a seeded 10/2 split of base-only rows, all other models use the
topology-held-out validation ID, the data-only model is trained only on labels
but selected with the common physics-aware validation score, and all 30 epochs
run without early stopping.}

\item \textbf{Topology selection and scenario generation were not
reproducible enough.}
\response{The manuscript now states the exact pre-split rule: retain the base
case plus the first 11 non-islanding outages in pandapower line-index order,
then shuffle outage IDs with NumPy \texttt{default\_rng(2026)} and floor the
0.70/0.15/0.15 fractions. No random or severity-based selection preceded this
cap. We explicitly limit representativeness because the selected set is not
severity-stratified and no descriptors were produced for excluded outages.
The graph-kernel length scales and the full synthetic diurnal construction
(clear-sky sinusoid, Beta cloud factor, Gaussian site factor, clipping) are
also specified.}

\item \textbf{Training and DC-OPF settings were incomplete.}
\response{The revised setup reports AdamW learning rates, weight decay,
gradient clipping, the 30-epoch cosine schedule, adaptive-weight Adam settings,
dropout, edge-conditioner dimensions, and the absence of warm restarts and
early stopping. DC-OPF is identified as pandapower 3.5.4
\texttt{rundcopp}, default PYPOWER backend, and $10^{-10}$ power tolerance,
with its active-power/angle model and omitted voltage/reactive variables stated
explicitly.}

\item \textbf{Five-seed bootstrap intervals were overinterpreted.}
\response{They are renamed descriptive seed-resampling intervals and are not
called confidence intervals. The revised table prints all five paired effects
for every case. The text reports that all 20 observed effects are positive
while explicitly declining population-level coverage or robustness claims.
The adaptive-weight figure now uses distinct line styles and markers so the
economic curve remains identifiable.}

\item \textbf{Latency and data-availability wording required correction.}
\response{``Times lower'' is replaced by an
$8.3$--$138.1\times$ raw device-specific latency ratio (speedup), with the
existing hardware and omitted-verification caveats retained. The availability
section names the exact ZIP and companion checksum JSON and states honestly
that a permanent repository DOI has not yet been assigned; it must be added
before publication.}
\end{enumerate}

The verification suite passes all 25 tests after these changes.

\end{document}
